% !TeX spellcheck = fr_FR
% --- --- --- --- --- --- --- --- --- --- --- --- --- --- --- --- %

\documentclass[a4paper,12pt]{report}  % args : article, report, book 
\usepackage[utf8]{inputenc}  % encodage UTF-8 
\usepackage[french]{babel}  % choix de langue, important pour maketitle et tableofcontents
\usepackage{geometry}  % géométrie de la page 
\usepackage{hyperref}  % liens internes et externes 
\usepackage{graphicx}  % images 
\usepackage{longtable}  % tableaux 
\usepackage{fontspec}  % polices d'écriture du système 
\usepackage{tabularx}  % pour des tableaux larges 
\usepackage{lipsum}  % pour du texte 
%\usepackage{titlepic}  % pour avoir une image dans \maketitle
\usepackage{titling}  % pour utiliser \theauthor, \thedate, \thetitle 
\usepackage{datetime}  % pour que la date de \thedate soit en français 

% --- --- --- --- --- --- --- --- --- --- --- --- --- --- --- --- %

\geometry{margin = 2 cm}  % marge de 2 cm (possible d'utiliser "in" aussi)

% hyperref 
\hypersetup{  
	colorlinks=true,
	linkcolor=black,
	urlcolor=blue,
	citecolor=red
}

% fontspec
\setmainfont{Bricolage Grotesque}
%\setmonofont[size=6pt]{Comic Code Ligatures} 


% --- --- --- --- --- --- --- --- --- --- --- --- --- --- --- --- %

\title{Cahier des charges fonctionnel}
\author{Groupe LRSVZZ}
\date{\today}
%\titlepic{\includegraphics[width=0.5\textwidth]{../Design/FichiersPoster/Logo_ESIEE.pdf}}

\begin{document}
	
%	\maketitle

	\begin{titlepage}
		\centering
		\vspace*{9cm} 
		{\LARGE \bfseries \thetitle \par}
		\vspace{1cm} 
		\large \theauthor \par
		\vspace{0.5cm} 
		\large \thedate \par
		\vspace{5cm} 
		\includegraphics[width=0.5\textwidth]{../Design/FichiersPoster/Logo_ESIEE.pdf} 
		\vfill 
	\end{titlepage}
	
%	\pagebreak
	
	% --- --- --- --- %
	
	\section*{Auteurs}
	
	\renewcommand{\arraystretch}{1.5}  % augmenter de 50 % les lignes 
	\begin{table}[h]
%		\centering
%		\caption{Auteurs}
%		\begin{tabular}{|c|c|c|}
		\begin{tabularx}{\textwidth}{|X|X|X|}  % on met X pour des colonnes de même longueur 
			\hline
			\textbf{Nom} & \textbf{Rôle} & \textbf{Département} \\
			\hline
			VARTANIAN Djivan & Responsable & Informatique \\
			\hline
			Row 2, Col 1 & Row 2, Col 2 & Row 2, Col 3 \\
			\hline
			Row 3, Col 1 & Row 3, Col 2 & Row 3, Col 3 \\
			\hline
%		\end{tabular}
		\end{tabularx}
	\end{table}
	
	\section*{Historique des documents}
	\begin{table}[h]
%		\centering
		\begin{tabular}{|c|c|c|c|} 
			\hline
			\textbf{Date} & \textbf{Version} & \textbf{Description du document} & \textbf{Auteur du document} \\
			\hline
			& & & \\
			\hline
			& & & \\
			\hline
			& & & \\
			\hline
		\end{tabular}
	\end{table}
	
	\section*{Approbations}
	\begin{table}[h]
		\begin{tabular}{|c|c|c|c|} 
			\hline
			\textbf{Date d'approbation} & \textbf{Version approuvée} & \textbf{Rôle de l'approbateur} & \textbf{Approbateur} \\
			\hline
			& & & \\
			\hline
			& & & \\
			\hline
			& & & \\
			\hline
		\end{tabular}
	\end{table}
	
	% --- --- --- --- %
		
	\tableofcontents
	
	\pagebreak
	% --- --- --- --- %
	
	\section{Résumé}
	% Brève description du projet, de ses objectifs et de ses résultats attendus.
	
	Résumé. 
	
	\begin{figure}[h]
		\centering
		\includegraphics[width=0.8\textwidth]{./attachements/Untitled-2025-06-11-1251_B.pdf}
		\caption{Description. }
		\label{fig:your_label}
	\end{figure}
	
	\section{Introduction}
	% Contexte du projet, problématique, objectifs du robot.
	Livrer des choses. Facilement, fiable. 
	
	\section{État de l'Art}
	% Revue des technologies existantes pour la livraison autonome.
	
	\section{Cahier des Charges Fonctionnel}
	\subsection{Description des Fonctionnalités}
	% Description des fonctionnalités requises (navigation, évitement d'obstacles, etc.)
	
	\subsection{Contraintes Techniques}
	% Contraintes techniques (dimensions, poids, autonomie, etc.)
	
	\section{Architecture du Système}
	\subsection{Schéma de l'Architecture}
	% Inclure un schéma de l'architecture matérielle et logicielle.
	
	\subsection{Description des Composants}
	% Description des composants utilisés (Raspberry Pi, capteur ultrason, etc.)
	
	\section{Développement et Implémentation}
	% Étapes de développement, utilisation de Git, défis rencontrés.
	
	\section{Tests et Validation}
	% Méthodologie de test, résultats des tests, retours d'expérience.
	
	\section{Rétro-planning et Gestion de Projet}
	\subsection{Organisation des Tâches}
	% Inclure un tableau ou un graphique du rétroplanning.
	
	\subsection{Liste des Actions Réalisées}
	% Liste des actions réalisées pendant le projet.
	
	\subsection{Compte Rendu des Réunions Hebdomadaires}
	% Résumé des points discutés lors des réunions.
	
	\section{Documentation et Ressources}
	% Documentation hébergée sur le serveur, liens vers les ressources.
	
	\section{Conclusion}
	% Bilan du projet, perspectives d'amélioration.
	
	\section{Annexes}
	% Graphiques, schémas, photos du robot, code source.
	
	\section{Remerciements}
	% Remerciements aux encadrants, partenaires, etc.
	
	
	
%	\section{Références}
%%	Listez toutes les références à des documents externes.
%	\lipsum[27]
%	
%	\section{Questions ouvertes}
%	\lipsum[28]
%	
%	\begin{longtable}{|l|l|l|l|l|l|l|l|}
%		\hline
%		\textbf{ID de la question} & \textbf{Enjeu} & \textbf{Élevé par} & \textbf{Décision} & \textbf{Résolu par} & \textbf{Statut} \\
%		\hline
%	\end{longtable}
%	
%	\section{Annexe}
%%	Inclure toute information supplémentaire ou documents annexes.
%	\lipsum[29]
	
\end{document}
